\section{The literal parity kernel}
\label{sec:tpc61-parity-kernel}

The normalized Gram in \cref{sec:tpc61-restricted-gram} is exact, but its
off-diagonal operator still appears there as an unequal-prime alias.  We now
calibrate every sign in that operator.  The result is a weighted correlation
of the two affine values
\[
 mj+h_0\quad\hbox{and}\quad m'j+h_0
\]
on the actual missing carrier.  No randomness hypothesis for the M\"obius
function is made, and no estimate for the resulting operator norm is asserted.

\subsection{Row gauge and exact sign calibration}

Retain the active row space \(S_M=\Ran\mathbf D_M\) and the unequal-prime
alias operator \(\Xi_{\ne p}\) from
\eqref{eq:tpc61-unequal-prime-operator}.  Let
\begin{equation}
 U_d:S_M\longrightarrow S_M,
 \qquad
 (U_dz)_m:=\mu(d_m)z_m.
 \label{eq:tpc61-source-mobius-gauge}
\end{equation}
The squarefree support and \((d_m,h_0)=1\) make \(U_d\) a diagonal
self-adjoint unitary.  Since \(\mathbf D_M\) is diagonal in the same row
basis, \(U_d\) commutes with every power and Moore--Penrose power of
\(\mathbf D_M\).

For two active rows define the literal common-output form
\begin{equation}
 \mathfrak C_{m,m'}
 :=
 \sum_{\substack{w,w'\in\cW_M\\
                   m(w)=m,\ m(w')=m'\\
                   i(w)=i(w')}}
 \mu(s(w))\mu(s(w'))
 \overline{\beta_w}\beta_{w'}.
 \label{eq:tpc61-common-output-form}
\end{equation}
For \(m=m'\), the fixed-\((i,m)\) singleton in
\eqref{eq:tpc61-ip-im-singletons} gives
\(\mathfrak C_{m,m}=W_M(m)\).  For \(m\ne m'\), the same expression is
the off-diagonal entry of \(M_X^*M_X\).

\begin{proposition}[Exact row-gauge calibration]
\label{prop:tpc61-exact-row-gauge}
Let \(m\ne m'\) be active missing rows.  Then
\begin{equation}
 \boxed{
 \bigl(U_d\Xi_{\ne p}U_d\bigr)_{m,m'}
 =\frac{1}{\sqrt{W_M(m)W_M(m')}}
 \sum_{\substack{w,w'\in\cW_M\\
                   m(w)=m,\ m(w')=m'\\
                   i(w)=i(w')}}
 \mu(n(w))\mu(n(w'))
 \overline{\beta_w}\beta_{w'}.}
 \label{eq:tpc61-two-affine-abstract-kernel}
\end{equation}
The diagonal of \(U_d\Xi_{\ne p}U_d\) is zero, and
\begin{equation}
 \lambda_{\max}(\Xi_{\ne p})
 =\lambda_{\max}(U_d\Xi_{\ne p}U_d).
 \label{eq:tpc61-row-gauge-spectrum}
\end{equation}
Consequently
\begin{equation}
 \kappa_M
 =1+\lambda_{\max}(U_d\Xi_{\ne p}U_d).
 \label{eq:tpc61-kappa-parity-gauge}
\end{equation}
\end{proposition}

\begin{proof}
On every high atom,
\[
 n(w)=p(w)r(w)
\]
is squarefree, \(p(w)\) is prime, and \((p(w),r(w))=1\).  Hence
\begin{equation}
 \mu(r(w))=-\mu(n(w)),
 \qquad
 \mu(s(w))
 =\mu(d_{m(w)})\mu(r(w))
 =-\mu(d_{m(w)})\mu(n(w)).
 \label{eq:tpc61-high-sign-calibration}
\end{equation}
The two minus signs cancel in a Gram product.  Conjugating its row indices
by \(U_d\) removes the remaining fixed factors
\(\mu(d_m)\mu(d_{m'})\), leaving exactly
\eqref{eq:tpc61-two-affine-abstract-kernel}.  The normalization by
\(W_M(m)^{-1/2}\) is the action of
\(\mathbf D_M^{\dagger/2}\) on the active row basis.  The diagonal
statement was proved after \eqref{eq:tpc61-native-gram-prime-split}; it
also follows directly because one pair \((i,m)\) determines one target and
one very-high prime.  Finally, unitary conjugation preserves eigenvalues,
which gives \eqref{eq:tpc61-row-gauge-spectrum}; combine this with
\eqref{eq:tpc61-kappa-unequal-prime}.
\end{proof}

The source-divisor sign has therefore been removed exactly; it has not been
discarded statistically.  The remaining arithmetic sign is the product of
the M\"obius values of two actual affine targets.  This is the parity-bearing
part of the native alias Gram.

\subsection{The direct-cell two-affine formula}

On the displayed left direct cell, write a high atom as
\(w=(m,j,\gamma)\), extending its coefficient by zero when the declared
literal support fails.  The factor dictionary of TPC--53 gives the exact
coefficient-free separation
\begin{equation}
 \beta_{m,j,\gamma}
 =a_m(j)K_{m\gamma}q_\gamma(j),
 \label{eq:tpc61-direct-cell-factorization}
\end{equation}
where \(a_m(j)\) contains the source and terminal shapes,
\(q_\gamma(j)\) contains the opposite-row shape and its fixed output
phase, and \(K_{m\gamma}\in\{0,1\}\) is the literal static mask.  The
fixed residue/orbit phase may be placed in either factor; the product in
\eqref{eq:tpc61-direct-cell-factorization} is the literal coefficient.

Let
\begin{equation}
 \mathfrak m_m(j)
 :=\one_{\{(m,j,\gamma)\in\cW_M
                 \text{ for some }\gamma\}},
 \label{eq:tpc61-missing-time-indicator}
\end{equation}
where coefficient zeros caused only by \(K_{m\gamma}\) or
\(q_\gamma(j)\) remain in the factors rather than in the indicator.  In
the identical-support specialization,
\begin{equation}
 \mathfrak m_m(j)=1
 \quad\Longleftrightarrow\quad
 \Pplus(mj+h_0)>y,
 \qquad
 r_m(j)\notin\cI_m,
 \label{eq:tpc61-missing-time-cofactor-test}
\end{equation}
where \(r_m(j)\) is defined in
\eqref{eq:tpc61-realized-row-cofactor}, and all common high-support tests
are understood.  Outside that specialization
one must replace the last test by the full first-failure signature.

Put
\begin{equation}
 \Gamma_{m,m'}(j)
 :=\sum_\gamma
 K_{m\gamma}K_{m'\gamma}|q_\gamma(j)|^2.
 \label{eq:tpc61-two-row-output-overlap}
\end{equation}

\begin{corollary}[Literal two-affine M\"obius kernel]
\label{cor:tpc61-two-affine-mobius-kernel}
For two distinct active rows \(m,m'\),
\begin{equation}
 \boxed{
 \begin{aligned}
 \bigl(U_d\Xi_{\ne p}U_d\bigr)_{m,m'}
 ={}&\frac{1}{\sqrt{W_M(m)W_M(m')}}
 \sum_j
 \mathfrak m_m(j)\mathfrak m_{m'}(j)\\
 &\times
 \mu(mj+h_0)\mu(m'j+h_0)
 \overline{a_m(j)}a_{m'}(j)
 \Gamma_{m,m'}(j).
 \end{aligned}}
 \label{eq:tpc61-two-affine-mobius-kernel}
\end{equation}
The sum is over the one fixed physical orbit slice.  The right direct slice
has the symmetric formula, and the two sides are joined only by orthogonal
direct sum.
\end{corollary}

\begin{proof}
A common native output on the left is precisely a common pair
\((\gamma,j)\).  Substitute
\eqref{eq:tpc61-direct-cell-factorization} into
\eqref{eq:tpc61-two-affine-abstract-kernel}, sum first over \(\gamma\),
and use \eqref{eq:tpc61-two-row-output-overlap}.  The missing and high
tests depend on \(m,j\), not on the surviving opposite coordinate, in the
identical-support specialization.
\end{proof}

Equation \eqref{eq:tpc61-two-affine-mobius-kernel} is an identity, not a
Chowla estimate.  In particular, cancellation for a single pair of rows
would not by itself control the largest eigenvalue when the number of rows
grows polynomially with \(X\).  The required statement is a weighted
matrix estimate for the complete family of growing slopes \(m\asymp Q\).

\subsection{Affine gcd separation}

The two affine targets in
\eqref{eq:tpc61-two-affine-mobius-kernel} have a simple exact resultant.

\begin{lemma}[Gcd and very-high-prime separation]
\label{lem:tpc61-affine-gcd-separation}
Let \(m\ne m'\) and suppose both targets at the same physical time \(j\)
are active very-high values:
\begin{equation}
 n=mj+h_0=pr,
 \qquad
 n'=m'j+h_0=p'r',
 \qquad p,p'>y.
 \label{eq:tpc61-two-high-factorizations}
\end{equation}
Then
\begin{equation}
 \boxed{\gcd(n,n')\mid h_0(m-m').}
 \label{eq:tpc61-affine-gcd-resultant}
\end{equation}
Moreover
\begin{equation}
 p\ne p',
 \qquad p\nmid n',
 \qquad p'\nmid n,
 \label{eq:tpc61-no-cross-high-prime}
\end{equation}
and hence neither very-high prime divides the other target's cofactor.
\end{lemma}

\begin{proof}
Every common divisor of \(n\) and \(n'\) divides the integer linear
combination
\[
 m'n-mn'=h_0(m'-m),
\]
which proves \eqref{eq:tpc61-affine-gcd-resultant}.  If \(p\mid n'\),
then, because \(p\mid n\), one has
\(p\mid(m-m')j\).  The inequality \(p>J_+\) excludes \(p\mid j\), and
\(|m-m'|\le\Delta_M<p\) excludes \(p\mid m-m'\).  Thus \(p\nmid n'\).
Interchanging the two rows gives \(p'\nmid n\), and in particular
\(p\ne p'\).  Since \(r'\mid n'\) and \(r\mid n\), the final assertion
follows.
\end{proof}

The gcd restriction controls local common factors, but it does not assign a
sign to
\(\mu(n)\mu(n')\).  It therefore cannot be entered as a cancellation
estimate for \(\Xi_{\ne p}\).

\subsection{Determinant reindexing}

For fixed distinct rows \(m,m'\), let
\(\mathscr S_{m,m'}\) consist of quadruples
\((p,r,p',r')\) satisfying all of the following:
\begin{enumerate}[label=\textup{(\roman*)},leftmargin=3em]
 \item \(p,p'>y\) are primes and \(pr,p'r'\) lie on the declared
       squarefree high support;
 \item
       \begin{equation}
        m'pr-mp'r'=h_0(m'-m);
        \label{eq:tpc61-double-prime-determinant}
       \end{equation}
 \item the common integer
       \begin{equation}
        j=\frac{pr-h_0}{m}
         =\frac{p'r'-h_0}{m'}
        \label{eq:tpc61-determinant-common-time}
       \end{equation}
       belongs to the fixed physical orbit slice; and
 \item \(p=\Pplus(pr)\), \(p'=\Pplus(p'r')\), and both recovered
       cofactors satisfy the literal missing tests.
\end{enumerate}
All inherited residue, profile, and endpoint conditions remain in this
definition; the determinant equation does not enlarge the carrier.

\begin{proposition}[Exact determinant form of one parity entry]
\label{prop:tpc61-determinant-parity-entry}
The map
\begin{equation}
 j\longmapsto
 \bigl(p_m(j),r_m(j),p_{m'}(j),r_{m'}(j)\bigr)
 \label{eq:tpc61-time-to-determinant-solution}
\end{equation}
is a bijection from the jointly missing times in
\eqref{eq:tpc61-two-affine-mobius-kernel} onto
\(\mathscr S_{m,m'}\).  Consequently
\begin{equation}
 \boxed{
 \begin{aligned}
 \bigl(U_d\Xi_{\ne p}U_d\bigr)_{m,m'}
 ={}&\frac{1}{\sqrt{W_M(m)W_M(m')}}
 \sum_{(p,r,p',r')\in\mathscr S_{m,m'}}
 \mu(r)\mu(r')\\
 &\times
 \overline{a_m(j)}a_{m'}(j)
 \Gamma_{m,m'}(j),
 \end{aligned}}
 \label{eq:tpc61-determinant-parity-kernel}
\end{equation}
where \(j\) is given by
\eqref{eq:tpc61-determinant-common-time}.
\end{proposition}

\begin{proof}
For a jointly missing time, unique largest-prime factorization supplies
\((p,r)\) and \((p',r')\).  Multiplying the first equality in
\eqref{eq:tpc61-two-high-factorizations} by \(m'\), the second by \(m\),
and subtracting gives
\eqref{eq:tpc61-double-prime-determinant}.  Thus the displayed map lands
in \(\mathscr S_{m,m'}\), and it is injective because either product
\(pr\) recovers \(j\).

Conversely, conditions \textup{(ii)}--\textup{(iii)} imply both affine
equalities at the same physical time, while \textup{(i)} and \textup{(iv)}
give exactly the active high and missing
tests.  Hence every quadruple comes from one jointly missing time.  Finally,
squarefreeness gives
\[
 \mu(pr)\mu(p'r')
 =(-\mu(r))(-\mu(r'))
 =\mu(r)\mu(r'),
\]
so \eqref{eq:tpc61-two-affine-mobius-kernel} becomes
\eqref{eq:tpc61-determinant-parity-kernel}.
\end{proof}

This reformulation exposes a possible bilinear or dispersion route, but it
does not prove that the signed determinant sum is smaller than its absolute
majorant.  The primes \(p,p'\) are distinct, yet distinctness alone gives no
M\"obius cancellation in \(r,r'\).

\subsection{Why the opposite-row mask does not diagonalize the kernel}

For a fixed time \(j\), put
\begin{equation}
 B(j):=\sum_\gamma|q_\gamma(j)|^2
 \label{eq:tpc61-opposite-row-mass}
\end{equation}
and, when \(B(j)>0\), define
\begin{equation}
 \eta_m(j)
 :=\frac{1}{B(j)}
   \sum_{\gamma:K_{m\gamma}=0}|q_\gamma(j)|^2.
 \label{eq:tpc61-row-mask-defect}
\end{equation}
Set \(\eta_m(j)=0\) when \(B(j)=0\).

\begin{proposition}[Exact two-row mask overlap]
\label{prop:tpc61-two-row-mask-overlap}
For every \(m,m',j\),
\begin{equation}
 \boxed{
 \bigl(1-\eta_m(j)-\eta_{m'}(j)\bigr)_+B(j)
 \le\Gamma_{m,m'}(j)\le B(j).}
 \label{eq:tpc61-two-row-mask-overlap-bound}
\end{equation}
On a nondegenerate complete direct opposite-row cell to which the
TPC--50 weighted mask theorem applies uniformly,
\begin{equation}
 \sup_{m,j}\eta_m(j)=o(1)
 \quad\Longrightarrow\quad
 \Gamma_{m,m'}(j)=(1-o(1))B(j)
 \label{eq:tpc61-mask-overlap-asymptotic}
\end{equation}
uniformly for all active \(m,m',j\).
\end{proposition}

\begin{proof}
The product \(K_{m\gamma}K_{m'\gamma}\) vanishes only on the union of
the two forbidden opposite-row sets.  The weight of that union is at most
the sum of their weights, which gives the lower bound.  The upper bound is
immediate from \(0\le K_{m\gamma}K_{m'\gamma}\le1\).  The final statement
follows from the uniform one-row defect estimate of TPC--50 and the same
union bound \citep{WangTPC50}.
\end{proof}

Thus static-mask survival has zero positive-power cost on the declared
interior cell, but it points in the opposite direction from an orthogonality
argument: two source rows retain almost the same opposite-row mass.  The
mask does not supply the open spectral estimate for
\(\Xi_{\ne p}\).  Any such estimate must use the signed affine correlation
in \eqref{eq:tpc61-two-affine-mobius-kernel}, the signed determinant form
in \eqref{eq:tpc61-determinant-parity-kernel}, or additional literal
structure not used here.

\subsection{The exact open spectral target}

The preceding results identify, but do not estimate, the remaining
uniform quantity.  Assume \(S_M\ne\{0\}\); when \(S_M=\{0\}\), the
missing map vanishes and this gate is vacuous.  Put
\begin{equation}
 \rho_{\rm par}(X)
 :=\lambda_{\max}(\Xi_{\ne p})
 =\lambda_{\max}(U_d\Xi_{\ne p}U_d)\ge0.
 \label{eq:tpc61-parity-spectral-target}
\end{equation}
The nonnegativity follows because \(\Xi_{\ne p}\) is Hermitian with zero
trace.  By
\eqref{eq:tpc61-kappa-unequal-prime},
\begin{equation}
 \boxed{\kappa_M=1+\rho_{\rm par}(X).}
 \label{eq:tpc61-kappa-parity-target}
\end{equation}
For one distinguished vector, only the corresponding Rayleigh quotient in
\eqref{eq:tpc61-pointwise-rayleigh} is needed.  A bound for that one
quadratic form is strictly weaker than
\eqref{eq:tpc61-parity-spectral-target} and must remain labeled
distinguished.

Neither the gcd lemma, the determinant bijection, Brun--Titchmarsh on one
cofactor ladder, nor the static-mask overlap proves
\(\rho_{\rm par}(X)=X^{o(1)}\).  Treating
\(\mu(mj+h_0)\) as an independent random sign would add an unproved
hypothesis and is not part of the TPC interface.
